
\subsubsection{2.1 Reinforcement Learning from Human Feedback}

RLHF was introduced for language model alignment by Ziegler et al. (2019) and scaled to production systems with InstructGPT \citep{ouyang2022instructgpt}. The standard pipeline involves: (1) supervised fine-tuning on demonstration data, (2) training a reward model on human preference rankings, and (3) optimizing the policy via PPO against the reward model with KL regularization.

Constitutional AI \citep{bai2022constitutional} introduced an alternative using AI feedback (RLAIF) with explicit constitutional principles for harmlessness training. Direct Preference Optimization \citep{rafailov2023dpo} eliminates the reward model, optimizing directly on preference pairs.

\subsubsection{2.2 Multi-Objective Alignment}

Askell et al. (2021) formalized the HHH framework (Helpful, Harmless, Honest), analyzing trade-offs between objectives. Their work demonstrated that multi-objective optimization is feasible but focused on helpfulness-harmlessness rather than helpfulness-controllability.

Gao et al. (2022) characterized reward model overoptimization, showing that excessive optimization degrades true reward. This informs the use of KL regularization in multi-objective settings.

\subsubsection{2.3 Instruction Following and Controllability}

IFEval (Zhou et al., 2023a) introduced a benchmark for verifiable instruction following, measuring compliance with 25 constraint types (format, length, keywords, structure). This work extends IFEval's use from evaluation to training signal extraction.

AlpacaEval \citep{dubois2024alpacaeval} measures instruction-following quality via GPT-4 judgment against reference outputs. AlpacaEval LC (length-controlled) is used as the helpfulness metric.

\subsubsection{2.4 Bidirectional Alignment Framework}

Sun et al. (2024) survey 400+ papers to propose a bidirectional taxonomy distinguishing $AI\rightarrow Human$ (specification integration---models adopting human values) and $Human\rightarrow AI$ (agency preservation---humans retaining control). Their framework is theoretical; this work provides an initial empirical implementation.
